% ------------------------------Info--------------------------------
% Dieses Dokument ist als Vorlage f�r die studentischen Arbeiten
% am Institut f�r Thermodynamik gedacht. Es stellt eine einfache 
% M�glichkeit dar, die studentischen Arbeiten im Corporate Design
% der TU Braunschweig zu gestalten.
% Dazu wurden viele der notwendigen Einstellungen bereits in dieses
% Dokument eingearbeitet. Es steht dem Autor nat�rlich frei weitere
% Anpassungen, an seine individuellen kreativen Vorstellungen und 
% zum erfolgreichen Erstellen notwendigen Bed�rfnisse, vorzunehmen.
% Hierf�r sei an dieser Stelle auf die Website der Corporate Design
% (https://www.tu-braunschweig.de/presse/cd/toolbox) und die 
% Dokumentation zum Corporate Design in LaTeX hingewiesen.
%
% In diesem Sinne: Frohes Schreiben!
%-------------------------------------------------------------------

\documentclass[fontsize=11pt,
               paper=a4,	%Papierformat ist A4
               style=print, %[<print/printdev/screen>] Dokument ist f�r den Druck gedacht
               nexus, %Hausschrift der TU Braunschweig verwenden
               cmyk, % [<cmyk/rgb>]  Farben in CMYK Schema
               twoside=false, %Kein zweiseitiger Druck
               toc=bibliography ,%Literaturverzeichnis ohne Nummer im Inhaltsverzeichnis anzeigen
               blue % [<orange/blue/green/violet>] Farbschema f�r die Arbeit
              ]
              {tubsbook}
                       
%---------------------------------------------------
%---------------Allgemeine Angaben------------------
%---------------------------------------------------
\def \TITLE					{EinhornFechten}
\def \TYPE					{Masterarbeit} %[<Projektarbeit/Bachelorarbeit/Studienarbeit/Masterarbeit>]
\def \AUTHOR				{Name des Studenten}
\def \BETREUER				{Name des Betreuers}
\def \TITELBILD				{pictures/Titelbild.png}
\def \AUFGABENSTELLUNG		{Aufgabenstellung.pdf}
\def \TOPLINE				{} %[<\showtopline>] % Aktiviert die rote Trennlinie �ber dem Bild. 
										%Sollte nur verwendet werden, wenn das Bild keine farbliche Trennung bietet.
%---------------------------------------------------
%----------Ende der Allgemeinen Angaben-------------
%
%		Header-Datei
%

% Nummerierungstiefe von section/subsection/subsubsection
\setcounter{secnumdepth}{3} 

%Absatzeinzug einstellen
\setlength{\parindent}{1em}
\setlength{\parskip}{0pt}

% Deutsche Rechtschreibung und deutscher Zeichensatz
\usepackage{german,ngerman}
\usepackage[ansinew]{inputenc}  

% Aktuelles Datum ermitteln
\usepackage[ngerman]{datenumber}

% Erweiterung f�r ein deutsches Literaturverzeichnis
\bibliographystyle{abbrvdin}

% Erweiterte Mathematikbibliotheken
\usepackage{amsmath}
\usepackage{amssymb}

% Ma�einhaeiten-Darstellung verbessern
\usepackage{units}

% Extra Symbole f�r die Verwendung im Text (z.B. \textcelsius)
%(werden nicht aus der Nexus-Schriftart einbunden)
\usepackage{textcomp}

% Einbinden von externen PDF Dateien
\usepackage[final]{pdfpages}

% Zum einbinden von Grafiken  
\usepackage{graphicx}
\usepackage{subcaption}  %erlaubt das erstellen von subfigures
\usepackage{float} %erm�glicht das erstellen von "non-floating figures"

%Hyperref f�r Verlinkungen im Dokument und URL's
%hyphens f�r verbesserten Zeilenumbruch in url's
\PassOptionsToPackage{hyphens}{url}
\usepackage[	breaklinks={true},
 						colorlinks	={true},
 					  	linkcolor	={tubsBlack},
  						citecolor	={tubsBlack},
  						filecolor	={tubsBlack},
  						urlcolor	={tuSecondaryMedium100},						
 					]{hyperref} %Interne Links in Schwarz und externe URL's in bg Farbe
 					
%Zus�tzliche Buchstaben und Zahlen erlauben f�r Zeilenumbruch bei url's (verbessert den Blocksatz)				
\expandafter\def\expandafter\UrlBreaks\expandafter{\UrlBreaks%  save the current one
  \do\a\do\b\do\c\do\d\do\e\do\f\do\g\do\i\do\j%
  \do\k\do\l\do\m\do\n\do\o\do\q\do\r\do\s%
  \do\u\do\v\do\w\do\x\do\y\do\z\do\A\do\B\do\C\do\D%
  \do\E\do\F\do\G\do\H\do\I\do\J\do\K\do\L\do\M\do\N%
  \do\O\do\P\do\Q\do\R\do\S\do\T\do\U\do\V\do\W\do\X%
  \do\Y\do\Z\do\1\do\2\do\3\do\4\do\5\do\6\do\7\do\8\do\9} 	
  
% R�ckseiten-Elemente
\address{% Addresse
  Technische Universit�t Braunschweig\\
  Institut f�r Thermodynamik\\
  Hans-Sommer-Strasse 5\\
  38106 Braunschweig}
\backpageinfo{}				

%Makro zur Formatierung des Anhangs
\newcommand{\Anhang}{
\appendix
\chapter*{Anhang}
\addcontentsline{toc}{chapter}{Anhang}
\setcounter{chapter}{1}
\markboth{Anhang}{Anhang}
\label{Anhang}}

%Makro zur Formatierung von der Titelseite bis zum Inhaltsverzeichniss
\newcommand{\Anfang}{
\pagenumbering{Roman}
\begin{titlepage}
%Designhelper f�r Erstellung 
%\showdesignhelper

\showtubslogo[right]
\showlogo{\includegraphics{ift_4C.pdf}}
\TOPLINE

%Titelbild f�r die Arbeit
\begin{titlerow}[bgimage=\TITELBILD ,imagefit=cropped]{3}


\hfill\vfill
\end{titlerow}
\begin{titlerow}[bgcolor=tuSecondaryDark100, c]{2}

\begin{fitbox}{\textwidth}{100pt}
\raggedright
\color{tubsWhite}
\textbf{\TITLE}
\end{fitbox}
\addvspace{1,5em}
\begin{Large}
\raggedright
\color{tubsWhite}
\textbf{\TYPE}\vfill
\end{Large}

\end{titlerow}
\begin{titlerow}[bgcolor=tuSecondaryMedium80]{3}
\color{tubsWhite}


\begin{tabbing}
\hspace*{2.0cm}   \=  \hspace*{12.0cm}        \= \kill
\textbf{Autor:} \> \textbf{\AUTHOR} \\
\textbf{Betreuer:} \> \textbf{\BETREUER}\\
\textbf{Pr\"ufer:} \> \textbf{Prof. Dr.-Ing. J\"urgen K\"ohler}\\


 \> \textbf{\today}
\end{tabbing}
\hfill\vfill
\end{titlerow}
\hfill\vfill
\end{titlepage}

\includepdf[pages=-]{\AUFGABENSTELLUNG}
\addchap*{Eidesstattliche Erkl�rung}


\vspace*{5cm}


Hiermit erkl�re ich eidesstattlich, dass ich diese Arbeit eigenst�ndig angefertigt und keine anderen als die angegebenen Hilfsmittel verwendet habe.
\vspace{3cm}\\
Braunschweig den \today
\tableofcontents
\pagenumbering{arabic}
\setcounter{page}{1}}



%Makro zum einpassen des Titels in die Titelbox
\usepackage{environ}
\newdimen\fontdim
\newdimen\upperfontdim
\newdimen\lowerfontdim
\newif\ifmoreiterations
\fontdim12pt
\makeatletter
\NewEnviron{fitbox}[2]{% \begin{fitbox}{<width>}{<height>} stuff \end{fitbox}
  \def\buildbox{%
    \setbox0\vbox{\hbox{\minipage{#1}%
      \fontsize{\fontdim}{1.2\fontdim}%
      \selectfont%
      \stuff%
    \endminipage}}%
    \dimen@\ht0
    \advance\dimen@\dp0
  }
  \def\stuff{\BODY}% Store environment body
  \buildbox
  % Compute upper and lower bounds
  \ifdim\dimen@>#2
    \loop
      \fontdim.5\fontdim % Reduce font size by half
      \buildbox
    \ifdim\dimen@>#2 \repeat
    \lowerfontdim\fontdim
    \upperfontdim2\fontdim
    \fontdim1.5\fontdim
  \else
    \loop
      \fontdim2\fontdim % Double font size
      \buildbox
    \ifdim\dimen@<#2 \repeat
    \upperfontdim\fontdim
    \lowerfontdim.5\fontdim
    \fontdim.75\fontdim
  \fi
  % Now try to find the optimum size
  \loop
    %\message{Bounds: \the\lowerfontdim\space
    %         \the\fontdim\space \the\upperfontdim^^J}
    \buildbox
    \ifdim\dimen@>#2
      \moreiterationstrue
      \upperfontdim\fontdim
      \advance\fontdim\lowerfontdim
      \fontdim.5\fontdim
    \else
      \advance\dimen@-#2
      \ifdim\dimen@<10pt
        \lowerfontdim\fontdim
        \advance\fontdim\upperfontdim
        \fontdim.5\fontdim
        \dimen@\upperfontdim
        \advance\dimen@-\lowerfontdim
        \ifdim\dimen@<.2pt
          \moreiterationsfalse
        \else
          \moreiterationstrue
        \fi
      \else
        \moreiterationsfalse
      \fi
    \fi
  \ifmoreiterations \repeat
  \box0% Typeset content
}
\makeatother %Header zum Einbinden der Packages und zur Layout formatierung/konfiguration
\begin{document}
\Anfang

%---------------------------------------------------
%--------------Beginn eigener Kapitel---------------
%---------------------------------------------------
\chapter{Control Theory}
\label{c:Control}

The following chapter introduces the needed basics of control theory.
In the first section a short introduction to the concepts of Multiple Input Multiple Output Control is 
given and the motivating problem is stated.

The second section deals with the process models and the system identification.

Afterwards the controller design techniques used in this thesis are explained.
Advantages and Disadvantages are given and several examples e.g. the Wood-Berry Destillation Column and Rosebrooks function are used. \par

\section{Problem Description from a Control Theoretic Point of View} % (fold)
\label{sec:Control:Problem}

\section{Controller Design Techniques}
\label{sec:Control:Design}

The problem described in \ref{sec:Control:Problem} is a standard issue in Control Theory, see e.g. \cite{Lunze2014}.

% section section_name (end) 
%\chapter{Grundlagen}
\label{sec:Grundlagen}



\section{Wert und Einheit}
\label{sec:Unit}
Viele Einheiten lassen sich sch�ner darstellen mit dem Latex Befehl \verb|\unit[]{}| beziehungsweise \verb|\unitfrac[]{}{}|. Siehe den Vergleich: ohne 1~m oder mit \unit[1]{m} bzw. ohne 1~m/sec oder mit \unitfrac[1]{m}{sec}. Des weiteren kann zum Beispiel mit \verb|\textcelsius| ein \textcelsius{} Symbol eingef�gt werden.

\section{�berschrift}
\label{sec:Quelltext}
Text Text Text Text Text Text Text Text Text Text Text Text Text Text Text Text Text Text Text Text Text Text Text Text Text Text Text Text Text Text Text Text Text Text Text Text Text Text Text Text Text Text Text Text Text Text Text Text Text Text Text Text Text Text Text Text Text Text Text Text Text Text Text Text Text Text Text Text Text Text Text Text

\section{Abbildungen einbinden}
\label{sec:pictures}
Text Text Text Text Text Text Text Text Text Text Text Text Text Text Text Text Text Text Text Text Text Text Text Text Text Text Text Text Text Text Text Text Text Text Text Text Text Text Text Text Text Text Text Text Text Text Text Text Text Text Text Text Text Text Text Text Text Text Text Text Text Text Text Text Text Text Text Text Text Text Text Text Text Text Text Text Text Text 


\begin{figure}[htbp]
  \centering
  \includegraphics{pictures/empty.jpg}
  \caption{Einzelne Abbildung}
  \label{fig:singlepicture}
\end{figure}


Text Text Text Text Text Text Text Text Text Text Text Text Text Text Text Text Text Text Text Text Text Text Text Text Text Text Text Text Text Text Text Text Text Text Text Text Text Text Text Text Text Text Text Text Text Text Text Text Text Text Text Text Text Text Text Text Text Text Text Text Text Text Text Text Text Text Text Text Text Text Text Text Text Text Text Text Text Text Text Text Text Text Text Text Text Text Text Text Text Text Text Text Text Text Text Text Text Text Text Text Text Text Text Text Text Text Text Text Text Text Text Text Text Text Text Text Text Text Text Text Text Text Text Text Text Text Text Text Text Text Text Text Text Text Text Text Text Text Text Text Text Text Text Text Text Text Text Text Text 



Text Text Text Text Text Text Text Text Text Text Text Text Text Text Text Text Text Text Text Text Text Text Text Text Text Text Text Text Text Text Text Text Text Text Text Text Text Text Text Text Text Text Text Text Text Text Text Text Text Text Text Text Text Text Text Text Text Text Text Text Text Text Text Text Text Text Text Text Text Text Text Text Text Text Text Text Text Text Text Text Text Text Text Text Text Text Text Text Text Text Text Text Text Text Text Text Text Text Text Text Text Text Text Text Text Text Text Text Text Text Text Text Text Text Text Text Text Text Text Text Text Text Text Text Text Text Text Text Text Text Text Text Text 
%\addchap{Symbolverzeichnis}
\label{Symbolverzeichnis}

\begin{tabbing}
\hspace*{3.5cm}                              \=  \hspace*{10.0cm}                     \= \kill
\textbf{Formelzeichen}                       \> \textbf{Bedeutung}                    \> \textbf{Einheit} \\
~\\
${I}$                                        \> Strom                                 \> A  \\
${i}$                                        \> Getriebe�bersetzung                   \> -  \\
${l}$                                        \> L�nge                                 \> m  \\
${M}$                                        \> Drehmoment                            \> N  \\
${m}$                                        \> Masse                                 \> Kg \\
${n}$                                        \> Drehzahl                              \> 1/sec \\
${R}$                                        \> Widerstand                            \> $\Omega$ \\
${t}$                                        \> Zeit                                  \> sec \\
${U}$                                        \> Spannung                              \> V \\
~\\
${\alpha}$                           		 \> Drehwinkel von Gelenk 1               \> rad \\
${\beta}$                           	 	 \> Drehwinkel von Gelenk 2               \> rad \\
${\Delta \varepsilon}$   				 	 \> Aufl�sung der Drehwinkelsensoren      \> rad/Imp. \\
${\varepsilon}$                      		 \> Drehwinkel allgemein                  \> rad \\
${\theta}$              					 \> Tr�gheitsmoment                       \> Kg$\cdot$m$^2$ \\
~\\
~\\
\textbf{Indizes}                             \> \textbf{Bedeutung} \\
~\\
${0}$                                        \> Leerlauf \\
${el}$                                       \> elektrischer Teil der Motoren \\
${mech}$                                     \> mechanischer Teil der Motoren \\
${G}$                                        \> Getriebe \\
${i, j}$                                	 \> Laufindex (1, 2, 3, ...) \\
${Ref}$                                      \> Referenzwert \\
${soll}$                                     \> Sollgr��e \\
\end{tabbing}

\bibliography{chapters/bibliography} 

\Anhang
\section{Erster Anhang}
\label{sec:Anhang_1}
Ein Anhang.


\section{Zweiter Anhang}
\label{sec:Anhang_2}
Ein weiterer Anhang.



%--------------------------------------------------
%--------------Ende eigener Kapitel----------------
%--------------------------------------------------

\makebackpage[trisec]%[<plain/info/addressinfo>]
\end{document}
